El análisis del flujo electromagnético permitió relacionar la disminución de una señal con la energía transportada y disipada en el material. Esta relación resultó pertinente para estudiar antenas inmersas, donde la atenuación del medio y la distribución del campo influyeron simultáneamente en la potencia observada.

La simulación FDTD permitió resolver los campos de una geometría finita y contrastarlos con la referencia de onda plana. Se empleó Meep por su formulación documentada, su interfaz de Python y su soporte de coordenadas cilíndricas \citep{Yee1966,Oskooi2010}. La simetría axial redujo el costo respecto de representar explícitamente la dirección azimutal; no se midió una aceleración frente a una simulación cartesiana 3D.

La contribución consistió en integrar la estimación de atenuación, sus controles de ajuste y una comprobación independiente del balance local. Se compararon los medios a una resolución común y se cuantificaron las limitaciones mediante refinamiento, variación temporal y registros de recursos. La conservación de archivos y versiones permitió revisar el procedimiento y reproducir los cálculos derivados.

La utilidad para telecomunicaciones se interpretó como una base física y numérica para estudios posteriores. Los resultados no constituyeron una caracterización de materiales medidos, una certificación de exposición ni un modelo validado de cobertura 5G. La selección de parámetros homogéneos a una frecuencia y el escalado geométrico delimitaron las conclusiones.
